% Options for packages loaded elsewhere
\PassOptionsToPackage{unicode}{hyperref}
\PassOptionsToPackage{hyphens}{url}
\documentclass[
]{article}
\usepackage{xcolor}
\usepackage{amsmath,amssymb}
\setcounter{secnumdepth}{-\maxdimen} % remove section numbering
\usepackage{iftex}
\ifPDFTeX
  \usepackage[T1]{fontenc}
  \usepackage[utf8]{inputenc}
  \usepackage{textcomp} % provide euro and other symbols
\else % if luatex or xetex
  \usepackage{unicode-math} % this also loads fontspec
  \defaultfontfeatures{Scale=MatchLowercase}
  \defaultfontfeatures[\rmfamily]{Ligatures=TeX,Scale=1}
\fi
\usepackage{lmodern}
\ifPDFTeX\else
  % xetex/luatex font selection
\fi
% Use upquote if available, for straight quotes in verbatim environments
\IfFileExists{upquote.sty}{\usepackage{upquote}}{}
\IfFileExists{microtype.sty}{% use microtype if available
  \usepackage[]{microtype}
  \UseMicrotypeSet[protrusion]{basicmath} % disable protrusion for tt fonts
}{}
\makeatletter
\@ifundefined{KOMAClassName}{% if non-KOMA class
  \IfFileExists{parskip.sty}{%
    \usepackage{parskip}
  }{% else
    \setlength{\parindent}{0pt}
    \setlength{\parskip}{6pt plus 2pt minus 1pt}}
}{% if KOMA class
  \KOMAoptions{parskip=half}}
\makeatother
\setlength{\emergencystretch}{3em} % prevent overfull lines
\providecommand{\tightlist}{%
  \setlength{\itemsep}{0pt}\setlength{\parskip}{0pt}}
\usepackage[]{natbib}
\bibliographystyle{plainnat}
\usepackage{bookmark}
\IfFileExists{xurl.sty}{\usepackage{xurl}}{} % add URL line breaks if available
\urlstyle{same}
\hypersetup{
  hidelinks,
  pdfcreator={LaTeX via pandoc}}

\author{}
\date{}

\begin{document}

Preprocessing was performed using \emph{QSIPrep} 26.0.0
\citep{cieslak2021qsiprep}, which is based on \emph{Nipype} 1.11.0
{[}\citet{nipype1}; \citet{nipype2}; RRID:SCR\_002502{]} and follows
\emph{NiPreps} framework \citep{nipreps}.

\paragraph{Anatomical data
preprocessing}\label{anatomical-data-preprocessing}

The T1-weighted (T1w) image was corrected for intensity non-uniformity
(INU) using \texttt{N4BiasFieldCorrection} \citep[ANTs 2.4.3]{n4}, and
used as an anatomical reference throughout the workflow. The anatomical
reference image was reoriented into AC-PC alignment via a 6-DOF
transform extracted from a full Affine registration to the
MNI152NLin2009cAsym template. A full nonlinear registration to the
template from AC-PC space was estimated via symmetric nonlinear
registration (SyN) using antsRegistration (\citet{ants}). Brain
extraction was performed on the T1w image using SynthStrip
\citep{synthstrip} and automated segmentation was performed using
SynthSeg \citep{synthseg1, synthseg2} from FreeSurfer version 7.3.1.

\paragraph{Diffusion data
preprocessing}\label{diffusion-data-preprocessing}

DWI data were denoised using the Marchenko-Pastur PCA method implemented
in dwidenoise
\citep{mrtrix3, dwidenoise1, dwidenoise2, cordero2019complex} with an
automatically-determined window size of 5 voxels. Any images with a
b-value less than 100 s/mm\^{}2 were treated as a \emph{b}=0 image. The
mean intensity of the DWI series was adjusted so all the mean intensity
of the b=0 images matched across eachseparate DWI scanning sequence. B1
field inhomogeneity was corrected using \texttt{dwibiascorrect} from
MRtrix3 with the N4 algorithm \citep{n4} after corrected images were
resampled.

FSL (version None)'s eddy was used for head motion correction and Eddy
current correction \citep{anderssoneddy}. Eddy was configured with a
\(q\)-space smoothing factor of 10, a total of 5 iterations, and 1000
voxels used to estimate hyperparameters. A quadratic first level model
and a linear second level model were used to characterize Eddy
current-related spatial distortion. \(q\)-space coordinates were
forcefully assigned to shells. Field offset was attempted to be
separated from subject movement. Shells were aligned post-eddy. Eddy's
outlier replacement was run \citep{eddyrepol}. Data were grouped by
slice, only including values from slices determined to contain at least
250 intracerebral voxels. Groups deviating by more than 4 standard
deviations from the prediction had their data replaced with imputed
values. Final interpolation was performed using the \texttt{jac} method.

Several confounding time-series were calculated based on the
preprocessed DWI: framewise displacement (FD) using the implementation
in \emph{Nipype} \citep[following the definitions by][]{power_fd_dvars}.
The head-motion estimates calculated in the correction step were also
placed within the corresponding confounds file. Slicewise cross
correlation was also calculated. The DWI time-series were resampled to
ACPC, generating a \emph{preprocessed DWI run in ACPC space} with 2mm
isotropic voxels.

Many internal operations of \emph{QSIPrep} use \emph{Nilearn} 0.10.1
\citep[RRID:SCR\_001362]{nilearn} and \emph{Dipy} \citep{dipy}. For more
details of the pipeline, see
\href{https://qsiprep.readthedocs.io/en/latest/workflows.html}{the
section corresponding to workflows in \emph{QSIPrep}'s documentation}.

\subsubsection{References}\label{references}

\bibliography{/out/logs/CITATION.bib}

\end{document}
